% section: 用語解説

\section{用語解説}\label{節：用語解説}
  本編や付録では,高校数学では直接習わないような用語も登場する.
  ここでは,漸化式を読むために役立つ用語を具体例から説明する.
  一度にすべて覚える必要はなく,意味が分からなくなったときに立ち戻ればよい.

  \subsection{線形・非線形}\label{節：用語解説：線形・非線形}
    まず,関数$f$が\textbf{線形}であるとは,入力の和と定数倍に対して
    \[
      f(x+y)=f(x)+f(y),\qquad
      f(cx)=cf(x)
    \]
    が成り立つことをいう.ここでは実数の入力と実数の定数$c$を考える.
    例えば$f(x)=2x$は線形である.
    一方,$f(x)=x^2$は一般に$(x+y)^2\ne x^2+y^2$なので線形ではない.
    このように,線形の条件を満たさないものを\textbf{非線形}という.

    漸化式では,未知の数列の項が係数を掛けて足す形だけで現れるものを\textbf{線形漸化式}という.
    例えば
    \begin{align*}
      a_{n+1}&=2a_n+3,\\
      a_{n+2}&=3a_{n+1}-2a_n
    \end{align*}
    は線形漸化式である.
    これに対して
    \[
      a_{n+1}=a_n^2+1,\qquad
      a_{n+1}=\frac{1}{a_n}
    \]
    は非線形漸化式である.後者では$a_n\ne0$が必要である.
    項の2乗,項どうしの積,項の逆数などが現れるかが見分ける目安になる.

    注意したいのは,線形かどうかは\textbf{何を未知の量として見ているか}で決まることである.
    \[
      a_{n+1}=n^2a_n+2^n
    \]
    も,$a_n$に関しては線形である.
    $n^2$や$2^n$は,\,$n$が決まれば既知の数だからである.
    一般項のグラフが直線になるという意味ではない.
    実際,$a_{n+1}=2a_n$は線形漸化式だが,その一般項は$a_n=a_1 2^{n-1}$である.

  \subsection{アフィン・同次・非同次}\label{節：用語解説：アフィン・同次・非同次}
    $f(x)=px+q$のように,線形な関数に定数を加えたものを\textbf{アフィン}な関数という.
    $q\ne0$なら$f(0)\ne0$であり,関数としては線形ではない.
    高校数学の一次関数は,$p\ne0$のアフィンな関数に当たる.
    それでも$a_{n+1}=pa_n+q$を線形漸化式と呼ぶのは,未知の項について線形な式に,既知の項を加えた方程式だからである.

    線形漸化式で,未知の項をすべて左辺に集めて
    \[
      a_{n+1}-pa_n=q
    \]
    と書いたとき,右辺が$0$なら\textbf{同次},そうでなければ\textbf{非同次}という.
    「同次」は「斉次」とも書く.
    右辺は定数に限らず,$n$についての既知の式でもよい.
    例えば$a_{n+1}-2a_n=n$は非同次である.

    同次線形漸化式では,2つの解の数列$u_n,v_n$から,定数$C,D$を用いて作る$Cu_n+Dv_n$も解になる.
    この性質を\textbf{重ね合わせの原理}という.
    例えば$u_{n+1}=2u_n$,\,$v_{n+1}=2v_n$なら
    \[
      Cu_{n+1}+Dv_{n+1}=2(Cu_n+Dv_n)
    \]
    となる.ただし,同じ初期条件を満たすとは限らない.
    また,非同次の場合は一般にこの性質をそのまま使えない.

    非同次の式を同次の式に直す例も見ておこう.
    $a_{n+1}=2a_n+3$に対して$b_n=a_n+3$と置くと,
    \[
      b_{n+1}=a_{n+1}+3=2(a_n+3)=2b_n
    \]
    となる.こうした置き換えが,本編で扱う解法の基本になる.

  \subsection{定数係数・変数係数}\label{節：用語解説：定数係数・変数係数}
    線形漸化式で未知の項に掛かる数を\textbf{係数}という.
    係数が$n$によらず一定なら\textbf{定数係数},\,$n$によって変わるなら\textbf{変数係数}という.
    例えば
    \[
      a_{n+2}=3a_{n+1}-2a_n
    \]
    は定数係数であり,
    \[
      a_{n+1}=(n+1)a_n
    \]
    は変数係数である.
    $a_{n+1}=2a_n+n^2$では,未知の項の係数は一定なので定数係数である.
    右辺の$n^2$は,未知の項に掛かる係数ではなく非同次の項である.
    定数係数かどうかと,同次かどうかは別々に判断する.

  \subsection{階数と初期条件}\label{節：用語解説：階数と初期条件}
    漸化式の\textbf{階数}は,関係する項の添字の最大値と最小値の差をいう.
    例えば$a_{n+1}=2a_n$は\textbf{一階},
    \[
      a_{n+2}=a_{n+1}+a_n
    \]
    は\textbf{二階}の漸化式である.
    二階は「隣接3項間」に対応するが,$a_{n+2}=a_n$も二階である.
    また,$a_{n+1}=a_n^2$は非線形だが一階であり,階数と項のべきの大きさは異なる.

    数列を具体的に決めるために最初に与える値を\textbf{初期条件}という.
    次の項が直前の$k$項から一意に計算できる$k$階の漸化式では,通常$a_1,\ldots,a_k$を指定する.
    上の二階の例なら$a_1,a_2$から$a_3,a_4,\ldots$を順に求められる.
    分母が$0$になるなど,更新の途中で式が定義されなくなる場合には別途注意が必要である.

  \subsection{反復・不動点}\label{節：用語解説：反復・不動点}
    $a_{n+1}=f(a_n)$では,同じ関数$f$を繰り返し適用して数列を作る.
    この操作を\textbf{反復}という.
    \[
      a_2=f(a_1),\qquad a_3=f(f(a_1))
    \]
    となり,$f(f(x))$は$f(x)$の2乗ではない.

    $f(\alpha)=\alpha$を満たす値$\alpha$を\textbf{不動点}という.
    その値から出発すれば,更新しても値が変わらない.
    例えば$f(x)=2x+3$の不動点は
    \[
      \alpha=2\alpha+3\iff\alpha=-3
    \]
    である.$a_1=-3$ならすべての項が$-3$になる.
    一般の初項でも,不動点からのずれ$b_n=a_n-\alpha$を考えると
    \[
      b_{n+1}=2b_n
    \]
    と簡単になる.先ほどの$b_n=a_n+3$という置き換えは,不動点からのずれを見ていたのである.
    不動点があっても,そこへ収束するとは限らない.
    この例では$a_1\ne-3$なら$|b_n|$は一歩ごとに2倍になる.

  \subsection{初等関数・閉じた形}\label{節：用語解説：初等関数・閉じた形}
    多項式,有理関数,指数関数,対数関数,三角関数,逆三角関数などを,有限回の四則演算や合成で組み合わせて得られる関数を\textbf{初等関数}という.
    合成とは,例えば$\log(1+x^2)$のように,ある関数の値を別の関数に入力することである.
    各関数や演算が定義される範囲で考える.

    一般項を\textbf{閉じた形}で表すという言い方もある.
    何を使ってよいかは文脈によるが,本書では初等関数による表現を一つの目安にしている.
    例えば$a_1=0$,\,$a_{n+1}=a_n+n$から
    \[
      a_n=\sum_{k=1}^{n-1}k=\frac{n(n-1)}{2}
    \]
    が得られる.最後の式なら,$n$が大きくなっても足し算の回数を増やさずに値を求められる.
    一方,$\sum_{k=1}^{n-1}k$も一般項の表現ではあるが,和に含まれる項数は$n$とともに増える.
    閉じた形が見つからなくても,項を順に計算したり,数列の性質を調べたりすることはできる.
